\documentclass[11pt,letterpaper]{article}
\usepackage[T1]{fontenc}
\usepackage[dvipsnames]{xcolor}
\usepackage{amsmath,amsthm,mathtools}
\usepackage{newpxtext,newpxmath}
\usepackage[margin=1in,headheight=14pt]{geometry}
\usepackage{microtype,booktabs,array,longtable,enumitem}
\usepackage{fancyhdr}
\usepackage[colorlinks=true,linkcolor=MidnightBlue,citecolor=MidnightBlue,urlcolor=MidnightBlue]{hyperref}
\usepackage[nameinlink,noabbrev]{cleveref}
\usepackage{bookmark}
\crefname{theorem}{Theorem}{Theorems}
\crefname{lemma}{Lemma}{Lemmas}
\crefname{proposition}{Proposition}{Propositions}
\crefname{corollary}{Corollary}{Corollaries}
\setlength{\parskip}{4pt}
\setlength{\parindent}{1.1em}
\setlist{itemsep=3pt,topsep=5pt}
\numberwithin{equation}{section}
\newtheorem{theorem}{Theorem}[section]
\newtheorem{lemma}[theorem]{Lemma}
\newtheorem{proposition}[theorem]{Proposition}
\newtheorem{corollary}[theorem]{Corollary}
\theoremstyle{definition}
\newtheorem{definition}[theorem]{Definition}
\newtheorem{conjecture}[theorem]{Conjecture}
\theoremstyle{remark}
\newtheorem{remark}[theorem]{Remark}
\DeclareMathOperator{\Li}{Li}
\DeclareMathOperator{\Imn}{Im}
\DeclareMathOperator{\supp}{supp}
\newcommand{\E}{\mathbb E}
\newcommand{\R}{\mathbb R}
\newcommand{\C}{\mathbb C}
\newcommand{\dd}{\,\mathrm d}
\newcommand{\ii}{\mathrm i}
\newcommand{\cut}{\C\setminus[1,\infty)}
\newcommand{\repo}{\texttt{ProveIt}}
\pagestyle{fancy}
\fancyhf{}
\fancyhead[L]{\small Global radial monotonicity}
\fancyhead[R]{\small ProveIt research continuation}
\fancyfoot[C]{\thepage}
\hypersetup{pdftitle={Global Radial Monotonicity for Harmonic Polylogarithms},pdfauthor={OpenAI-assisted research for the ProveIt project},pdfsubject={Exact all-radius inequalities, quantitative stability and positive binomial identities}}
\begin{document}
\begin{titlepage}
\setlength{\parindent}{0pt}
\thispagestyle{empty}
{\small\sffamily\bfseries PROVEIT\quad /\quad ANALYSIS\quad /\quad POLYLOGARITHMS\par}
\vspace{17mm}
{\Huge\bfseries Global Radial Monotonicity\par}
\vspace{3mm}
{\Huge\bfseries for Harmonic Polylogarithms\par}
\vspace{8mm}
{\Large A uniform outer-order theorem, quantitative\par
inner-order stability, and positive binomial identities\par}
\vspace{13mm}
{\large Research continuation prepared for\par
Vladimir Reshetnikov's ProveIt project\par}
\vspace{4mm}
{10 October 2026\par}
\vspace{12mm}
\begin{minipage}{\textwidth}
\textbf{Abstract.}
For $F_{a,b}(z)=\sum_{n>m\ge1}z^n/(n^a m^b)$, let
$\theta_{a,b}(\rho)$ be the unique upper-semicircle zero of its imaginary
part and put $\eta_{a,b}(\rho)=\cos\theta_{a,b}(\rho)/\rho$.
We prove strict decrease of $\eta$ on the entire interval $0<\rho\le1$
for every real $a\ge8$ and every $b>0$, with the explicit differential
bound $\eta'(\rho)<-\rho H_4^{(b)}(2/5)^a/120$.
The argument combines a uniform localization of the zero, a low-mode
inequality certified by ten positive rational Bernstein coefficients,
and analytic bounds for the full infinite tail.
A separate stability theorem proves the same global direction for every
$a>0$ whenever $2^{-b}<[(2/5)^a-(3/8)^a]/256$.
Thus the integer outer-order conjecture is reduced to six bounded
inner-order strips, rather than merely verified near the origin.
We derive a weighted cubic-moment criterion for radial motion, a
positive squared-resolvent representation, and convergent binomial
identities with exact remainders, including a specialization for
$\pi\log2$. Standard-library rational verifiers certify the uniform
bounds, eight angular-root brackets, and nine Gaussian values.
The full normalized-radius conjecture and the arithmetic $S_6,S_8$
candidates are not claimed resolved.
\end{minipage}
\vfill
{\small Ordinary analytic proofs with finite exact certificates; not a proof-assistant
formalization or an externally refereed paper. No literature-wide priority claim is made.\par}
\end{titlepage}
\begingroup
\small
\setlength{\parskip}{0pt}
\tableofcontents
\endgroup
\newpage

\section{Research target, provenance, and principal results}
\label{sec:overview}
The current manuscript \emph{Polylogarithms and their Arithmetic Bridges}
already proves that the first nonconstant radial coefficient has the
conjectured sign for all real outer orders $a\ge2$ and all positive
inner orders $b$ \cite{repo-local}. That is a local theorem: it does not
exclude a derivative changing sign farther along the zero curve.
The manuscript also proves global decrease in the exact limit
$b=\infty$, while explicitly declining to infer the same assertion for
finite $b$ \cite{repo-cartesian}. These are the two gaps addressed here.

Our repository reference is the retrieved revision
\begin{center}
\small\texttt{3338e535ef54ad8f6be56bdce8b3d084a2924465}.
\end{center}
The relevant chapter paths and blob identifiers are recorded in
\path{provenance/SOURCES.md}. The incoming directory was inventoried;
this is not a claim that every incoming archive was fully audited.
The accessible incoming quartic-bifurcation article treats a distinct,
local fractional-order transition and two turning points
\cite{incoming-bifurcation}. The incoming all-depth and sharp-Euler
articles likewise address different questions
\cite{incoming-oscillation,incoming-euler}.

\subsection{Conventions and inherited facts}
For $a,b>0$, write
\begin{equation}\label{eq:def}
 F_{a,b}(z)=\sum_{n\ge2}c_nz^n,
 \qquad c_n=\frac{H_{n-1}^{(b)}}{n^a},
 \qquad H_k^{(b)}=\sum_{j=1}^k j^{-b}.
\end{equation}
The sum is strict, not a multiple-star sum. We use the principal
continuation to $\cut$, consistent with the depth-one conventions in
DLMF \S25.12 \cite{dlmf-polylog}. Every boundary value below lies away
from the cut. In both parameter regimes of our global theorems,
$a+b>1$, so the boundary series is also ordinarily convergent.

The manuscript's finite signed-measure theorem supplies the following
fact on the domain $a,b>0$, $a+b>1$: on every open upper semicircle
$z=\rho e^{\ii\theta}$, $0<\rho\le1$, the imaginary part has exactly one
simple zero, denoted $\theta_{a,b}(\rho)$, with
\begin{equation}\label{eq:inherited-angle}
 \arccos\rho<\theta_{a,b}(\rho)<\frac\pi2,
 \qquad \partial_\theta\Imn F_{a,b}(\rho e^{\ii\theta})<0
 \quad\text{at the zero}.
\end{equation}
This statement and its proof are inherited from \cite{repo-orders},
not claimed as a new result. For $a\ge8$ we also give an independent
uniqueness argument using the same estimates as the radial proof.
Define
\begin{equation}\label{eq:normalization}
 s=\rho^2,\qquad
 m(s)=\frac{2\cos\theta_{a,b}(\rho)}\rho,\qquad
 \eta_{a,b}(\rho)=\frac{m(\rho^2)}2.
\end{equation}
In particular, $\eta'(\rho)=\rho m'(s)$; this factor is essential.

\subsection{Statements of the two global results}
\begin{theorem}[Uniform outer-order theorem]\label{thm:main}
For every real $a\ge8$, every $b>0$, and $0<\rho\le1$,
\begin{equation}\label{eq:main-bound}
 \eta_{a,b}'(\rho)
 <-\frac{\rho}{120}H_4^{(b)}\left(\frac25\right)^a<0.
\end{equation}
The function extends continuously to zero with
\begin{equation}\label{eq:initial}
 \eta_{a,b}(0)=\frac{1+2^{-b}}2\left(\frac23\right)^a,
\end{equation}
and consequently
\begin{equation}\label{eq:gap}
 \eta_{a,b}(0)-\eta_{a,b}(\rho)
 >\frac{\rho^2}{240}H_4^{(b)}\left(\frac25\right)^a.
\end{equation}
At $\rho=1$, the derivative is the derivative of the local analytic
continuation of the zero curve, equivalently its left derivative.
\end{theorem}

\begin{theorem}[Quantitative finite-inner-order transfer]\label{thm:large-b}
For $a>0$, put
\begin{equation}\label{eq:delta}
 \Delta(a)=\left(\frac25\right)^a-\left(\frac38\right)^a>0.
\end{equation}
If $b>0$ satisfies
\begin{equation}\label{eq:large-b-condition}
 2^{-b}<\frac{\Delta(a)}{256},
\end{equation}
then $\eta_{a,b}'(\rho)<0$ for every $0<\rho\le1$.
Under the stronger condition $2^{-b}\le\Delta(a)/512$,
\begin{equation}\label{eq:large-b-gap}
 \eta_{a,b}'(\rho)\le-\frac{\rho\Delta(a)}{16}.
\end{equation}
\end{theorem}

For integer $a$ the first sufficient integer thresholds in
\eqref{eq:large-b-condition} are as follows; these are not optimal
thresholds for the actual phenomenon.
\begin{center}
\begin{tabular}{@{}lrrrrrrrr@{}}
\toprule
Outer order $a$ &1&2&3&4&5&6&7&8\\
Sufficient $b\ge$&14&14&15&16&17&18&19&20\\
\bottomrule
\end{tabular}
\end{center}
For $a=8$ the entry is redundant by \cref{thm:main}. Thus the original
integer conjecture for $a\ge2$ is reduced to the six strips
$2\le a\le7$, with the integer $a$ fixed and $0<b<B_a$ bounded by this
table. No proof is claimed for every point of those remaining strips.

\subsection{Scope and originality discipline}
The additions relative to the inspected sources are the two explicit
all-radius regimes, their quantitative gaps, the cubic-moment velocity
formula, and the stated binomial identities and certificates.
Positive integral manipulations, Chebyshev representations, Bernstein
positivity, and elementary geometric acceleration are background
methods. An identity derived here need not be absent from the wider
literature. In particular, the new positive series is not asserted to
improve the sharp universal Euler evaluator in the incoming work.

\section{Positive kernels, moments, and a real zero equation}
\label{sec:kernels}
\subsection{The inherited double kernel and its probability form}
The manuscript proves \cite{repo-kernel}
\begin{equation}\label{eq:double}
 F_{a,b}(z)=z^2\int_{0<y<x<1}
 \frac{W_{a,b}(x,y)}{(1-xz)(1-yz)}\dd y\dd x,
\end{equation}
where
\begin{equation}\label{eq:weight}
 W_{a,b}(x,y)=
 \frac{(-\log x)^{a-1}(\log(x/y))^{b-1}}{\Gamma(a)\Gamma(b)},
 \qquad \int W_{a,b}=2^{-a}=c_2.
\end{equation}
For completeness, the substitution $y=xv$ gives
\[
 \int W_{a,b}(x,y)x^j y^k\dd y\dd x
   =\frac1{(j+k+2)^a(k+1)^b}.
\]
Geometric expansion and summation over $j+k=n-2$ recover $c_n$.
Uniform domination on compact subsets of $\cut$ proves continuation.
The positivity of this kernel will be used, but does not by itself
settle the sign of a normalized radial derivative.

Let $Y,Z$ be independent gamma variables with shapes $a,b$ and rates
$2,1$, respectively. Set $X=e^{-Y}$ and $V=e^{-Z}$. The same substitution
rewrites \eqref{eq:double} as
\begin{equation}\label{eq:probability}
 F_{a,b}(z)=c_2z^2\E\frac1{(1-Xz)(1-XVz)},
 \qquad \E X^k=\left(\frac2{k+2}\right)^a,
 \quad \E V^k=(k+1)^{-b}.
\end{equation}
All variables lie strictly between zero and one.

\subsection{A positive squared-resolvent measure}
\begin{proposition}\label[proposition]{prop:squared}
There is a positive finite measure $\mu=\mu_{a,b}$, with full support
$[0,1]$ and mass $c_2$, such that
\begin{equation}\label{eq:squared}
 F_{a,b}(z)=z^2\int_0^1\frac{\dd\mu(u)}{(1-uz)^2},
 \qquad
 \int_0^1u^k\dd\mu(u)=\frac{c_{k+2}}{k+1}\quad(k\ge0).
\end{equation}
In particular, for $\ell,k\ge0$,
\begin{equation}\label{eq:hausdorff}
 \sum_{j=0}^k(-1)^j\binom kj
 \frac{H_{\ell+j+1}^{(b)}}{(\ell+j+1)(\ell+j+2)^a}>0.
\end{equation}
Every finite Hankel moment matrix of $\mu$ is positive definite.
\end{proposition}
\begin{proof}
Use the elementary identity
\[
 \frac1{(1-xz)(1-yz)}=
 \int_0^1\frac{\dd t}{[1-(tx+(1-t)y)z]^2}
\]
in \eqref{eq:double}. The image of the positive product measure under
$u=tx+(1-t)y$ defines $\mu$. Every nonempty open subinterval of $(0,1)$
has positive mass, so the support is full. Comparing coefficients in
$(1-uz)^{-2}=\sum_{k\ge0}(k+1)u^kz^k$ proves the moments.
The left side of \eqref{eq:hausdorff} is
$\int u^\ell(1-u)^k\dd\mu>0$. For a nonzero polynomial $p$,
$\int p(u)^2\dd\mu(u)>0$, proving the Hankel assertion.
\end{proof}
This is a positive representation with a \emph{squared} resolvent. It
does not contradict the failure of a finite signed \emph{first}-resolvent
representation in part of the low-order domain.

\subsection{Two forms of the zero equation}
For $s=\rho^2$ and $m=2\cos\theta/\rho$ let
\begin{equation}\label{eq:D}
 D_u(m,s)=1+s(u^2-mu).
\end{equation}
Rationalizing \eqref{eq:squared} yields
\begin{equation}\label{eq:P-integral}
 \Imn F_{a,b}(\rho e^{\ii\theta})
 =\rho^3\sin\theta\,P(m,s),\qquad
 P(m,s)=\int\frac{m-2u}{D_u(m,s)^2}\dd\mu(u).
\end{equation}
Similarly \eqref{eq:probability} yields
\begin{equation}\label{eq:G-b}
 P(m,s)=c_2G_b(m,s),\qquad
 G_b(m,s)=\E\frac{m-X-XV}{D_X(m,s)D_{XV}(m,s)}.
\end{equation}
In particular $P(0,s)<0$ for all $0\le s\le1$, and the imaginary part
is negative at every upper-half-disk point with $\cos\theta\le0$.

For the series version, let $U_n$ denote the Chebyshev polynomial of the
second kind. Its standard trigonometric representation is
$U_{n-1}(\cos\theta)=\sin(n\theta)/\sin\theta$
\cite{dlmf-chebyshev}. Define
\begin{equation}\label{eq:T-def}
 T_n(m,s)=\rho^{n-3}U_{n-1}(\rho m/2).
\end{equation}
Parity shows that this is a polynomial in $m,s$ for $n\ge2$, and
\begin{equation}\label{eq:T-recurrence}
 T_2=m,\quad T_3=sm^2-1,\quad
 T_{n+1}=s(mT_n-T_{n-1}),\qquad
 P(m,s)=\sum_{n\ge2}c_nT_n(m,s).
\end{equation}
The next four polynomials are
\begin{align}\label{eq:T-low}
 T_4&=s^2m^3-2sm,&
 T_5&=s^3m^4-3s^2m^2+s,\notag\\
 T_6&=s^4m^5-4s^3m^3+3s^2m,&
 T_7&=s^5m^6-5s^4m^4+6s^3m^2-s^2.
\end{align}
At $s=0$ the equation is $c_2m-c_3=0$, giving
$m(0)=c_3/c_2$ and \eqref{eq:initial}.

\section{Uniform localization and control of the angular derivative}
\label{sec:localization}
From now until the end of \cref{sec:outer-proof}, assume $a\ge8$, $b>0$.
Write
\begin{equation}\label{eq:UVW}
 x=2^{-b},\quad y=3^{-b},\quad
 U=1+x,\quad V_0=1+x+y,\quad W=1+x+y+x^2.
\end{equation}
We distinguish $V_0$ from the random variable $V$ in
\eqref{eq:probability}. Then
$c_3=U3^{-a}$, $c_4=V_0 4^{-a}$, $c_5=W5^{-a}$, and
\begin{equation}\label{eq:xy}
 0<x<1,\qquad x^2\le y\le x,\qquad
 \frac{U^2}{V_0}\le\frac43.
\end{equation}
The last inequality follows from
$(1+x)^2/(1+x+x^2)\le4/3$.

\begin{lemma}[Uniform coefficient comparisons]\label[lemma]{lem:coefficients}
For the indicated ranges of $n$,
\begin{align}\label{eq:coeff-bounds}
 \frac{c_n}{c_2}&\le(n-1)\left(\frac2n\right)^a &&(n\ge2),\notag\\
 \frac{c_n}{c_3}&\le\frac{n-1}{2}\left(\frac3n\right)^a &&(n\ge3),\\
 \frac{c_n}{c_5}&\le\frac{n-1}{4}\left(\frac5n\right)^a &&(n\ge5).\notag
\end{align}
\end{lemma}
\begin{proof}
The averages $H_k^{(b)}/k$ decrease because the summands $j^{-b}$
decrease. Compare these averages at $k=n-1$ and at $k=1,2,4$.
\end{proof}

\begin{lemma}[Elementary Chebyshev bounds]\label[lemma]{lem:chebyshev}
For $n\ge3$ and $0\le c\le1/6$,
\begin{equation}\label{eq:Uprime}
 |U_{n-1}'(c)|\le\frac{11}{10}n.
\end{equation}
For $0\le c\le1/20$,
\begin{equation}\label{eq:U-parity}
 |U_{n-1}(c)|\le
 \begin{cases}
 101/100,&n\text{ odd},\\
 (101/100)nc,&n\text{ even}.
 \end{cases}
\end{equation}
For all $-1\le c\le1$, $|U_{n-1}(c)|\le n$.
\end{lemma}
\begin{proof}
Differentiating the trigonometric representation gives
\[
 |U_{n-1}'(c)|\le\frac n{1-c^2}
                 +\frac c{(1-c^2)^{3/2}}
 \le\frac{36n}{35}+\frac15<\frac{11n}{10}.
\]
Here $n\ge3$ is used in the final comparison.
Write $\theta=\pi/2-\delta$, so $\sin\delta=c$.
For odd $n$, use $|\sin(n\theta)|\le1$; for even $n$, use
$|\sin(n\delta)|\le n\sin\delta$. Division by
$\sin\theta=\sqrt{1-c^2}>100/101$ gives \eqref{eq:U-parity}.
Finally $|\sin(n\theta)|\le n\sin\theta$ proves the global bound,
including endpoints by continuity.
\end{proof}

Set
\begin{equation}\label{eq:C-M}
 C=\frac{11}{10},\qquad q=\frac{c_3}{c_2},\qquad
 M=Cq,\qquad M_8=\frac{11}{5}\left(\frac23\right)^8
 =\frac{2816}{32805}<\frac1{10}.
\end{equation}
Three rationally certified bounds used below are
\begin{align}\label{eq:certificate-three}
 B_8&:=\sum_{n\ge3}n(n-1)(2/n)^8<\frac13,\notag\\
 C\left(1-\frac{101}{200}E_8\right)-\frac{101}{100}O_8
 &>\frac1{50},\\
 O_8&:=\sum_{\substack{n\ge3\\ n\text{ odd}}}
          \frac{n-1}{2}(3/n)^8,\quad
 E_8:=\sum_{\substack{n\ge4\\ n\text{ even}}}n(n-1)(2/n)^8.\notag
\end{align}
The exact finite-sum and integral-tail verification is given in
\cref{app:certificates}; no decimal approximation is an acceptance test.

\begin{proposition}[Root strip and derivative bounds]\label[proposition]{prop:strip}
Every upper-semicircle zero lies in
\begin{equation}\label{eq:strip}
 0<m(s)<C\frac{c_3}{c_2}\le M_8<\frac1{10}.
\end{equation}
It is the only such zero, is simple, and on the larger strip
$0\le m\le1/3$, $0\le s\le1$,
\begin{equation}\label{eq:Pm}
 \frac45c_2<P_m(m,s)<\frac65c_2.
\end{equation}
\end{proposition}
\begin{proof}
There is no zero with $\cos\theta\le0$ by \eqref{eq:G-b}.
For any zero with $m>0$, the global Chebyshev bound and
\eqref{eq:coeff-bounds} give
\[
 c_2m\le\sum_{n\ge3}n c_n\rho^{n-3}\le c_2B_8,
\]
so all possible zeros lie in $m<1/3$.
On this strip $c=\rho m/2\le1/6$ and
$\partial_mT_n=\rho^{n-2}U_{n-1}'(c)/2$. Therefore
\[
 |P_m-c_2|\le\frac{11}{20}\sum_{n\ge3}n c_n
 \le\frac{11}{20}B_8c_2<\frac{11}{60}c_2,
\]
which implies \eqref{eq:Pm}.

We have $P(0,s)<0$. At $m=M$, the parity bounds apply because
$M\le M_8<1/10$. Bound all odd-indexed terms $n\ge3$ and all even-indexed
terms $n\ge4$ in absolute value, while retaining $c_2M$:
\[
 \frac{P(M,s)}{c_3}\ge
 C\left(1-\frac{101}{200}E_8\right)-\frac{101}{100}O_8
 >\frac1{50}.
\]
The coefficient comparisons and the inequalities $(r/n)^a\le(r/n)^8$
for $n\ge r$ justify the uniform replacements. Existence follows by
the intermediate value theorem, and strict increase of $P$ on the
containing strip proves uniqueness and simplicity.
The estimates also justify termwise differentiation: the relevant
majorants are bounded by constant multiples of $\sum n^{2-a}$.
At zero the implicit-function theorem applies to $P(m,0)=c_2m-c_3$.
Away from zero, analyticity of $F$ off the cut and simplicity give the
local analytic continuation, including through $s=1$.
\end{proof}

\section{The uniform radial derivative inequality}
\label{sec:outer-proof}
We now prove the main sign. The distinction between localization and
radial monotonicity matters: a small zero does not by itself fix the
sign of its velocity.

\subsection{A positive low-mode contribution}
\begin{lemma}[The first five relevant modes]\label[lemma]{lem:low}
At every point with $0\le m\le M$, $0\le s\le1$,
\begin{equation}\label{eq:low-derivative}
 \sum_{n=2}^6 c_n\partial_sT_n
 \ge c_5+c_3m^2-2c_4m-6c_5m^2
 >\frac25c_5.
\end{equation}
\end{lemma}
\begin{proof}
From \eqref{eq:T-low},
\[
 \partial_sT_3=m^2,\quad
 \partial_sT_4\ge-2m,\quad
 \partial_sT_5\ge1-6m^2,
\]
and
\[
 \partial_sT_6=sm(6-12sm^2+4s^2m^4)\ge0
\]
because $m<1/10$. This proves the first inequality.
Moreover
\[
 \frac{M}{c_4/c_3}
 =C\frac{U^2}{V_0}\left(\frac89\right)^a
 \le2C\left(\frac89\right)^8<1.
\]
Thus $c_3m^2-2c_4m$ decreases on $[0,M]$. Substitution of $M=Cc_3/c_2$
and $m^2\le M^2$ gives a lower bound $c_5\mathcal E(a,b)$, where
\begin{equation}\label{eq:core}
 \mathcal E(a,b)=1-2C\frac{UV_0}{W}\left(\frac56\right)^a
  +C^2\frac{U^3}{W}\left(\frac{20}{27}\right)^a
  -6C^2U^2\left(\frac49\right)^a.
\end{equation}
We prove that $\mathcal E(a,b)>2/5$ uniformly.

First, $\mathcal E(a,b)$ increases with $a\ge8$. For the two middle
exponentials, the ratio of the magnitude of the negative derivative to
the positive derivative is
\[
 \frac{CU^2}{2V_0}\left(\frac89\right)^a
       \frac{\log(27/20)}{\log(6/5)}
 <\frac{22}{15}\left(\frac89\right)^8<1.
\]
We used $U^2/V_0\le4/3$ and
$27/20<(6/5)^2$, which implies the logarithm ratio is less than two.
The last term in \eqref{eq:core} also has positive derivative in $a$.

It remains to verify $\mathcal E(8,b)>2/5$. Multiply the difference
by the positive $W$. The resulting polynomial is
\begin{equation}\label{eq:core-polynomial}
 p(x,y)=\frac35W-2C(5/6)^8UV_0
        +C^2(20/27)^8U^3-6C^2(4/9)^8U^2W.
\end{equation}
It is affine in $y$. On the entire larger region
$0\le x\le1$, $x^2\le y\le x$, positivity therefore follows from
positivity at $y=x^2$ and $y=x$. Each endpoint polynomial has degree at
most four. In the degree-four Bernstein basis
$\binom4j x^j(1-x)^{4-j}$, all ten coefficients are strictly positive.
Their exact numerators with a common denominator are in
\cref{tab:bernstein}. This proves positivity on the full region,
including the actual curve $y=3^{-b}$, and finishes the proof.
\end{proof}

\begin{table}[htbp]
\centering
\caption{Exact Bernstein certificate for \eqref{eq:core-polynomial}.
Divide every displayed integer by $7230196133913600$.}
\label{tab:bernstein}
\begin{tabular}{@{}crr@{}}
\toprule
$j$ & $p(x,x^2)$ & $p(x,x)$\\
\midrule
0&1351861016318884&1351861016318884\\
1&1121529834868342&1261248688948063\\
2&1434025919811540& 977124901564326\\
3&1462870427049383& 617821376075897\\
4& 221755156299224& 221755156299224\\
\bottomrule
\end{tabular}
\end{table}

\subsection{A bound for the entire remaining tail}
\begin{lemma}[Radial tail]\label[lemma]{lem:tail}
For $0\le m\le M_8$ and $0\le s\le1$,
\begin{equation}\label{eq:radial-tail}
 \left|\sum_{n\ge7}c_n\partial_sT_n(m,s)\right|
 <\frac{39}{100}c_5.
\end{equation}
\end{lemma}
\begin{proof}
For $s>0$, differentiation of \eqref{eq:T-def} gives
\[
 \partial_sT_n=
 \frac{n-3}{2}\rho^{n-5}U_{n-1}(\rho m/2)
 +\frac m4\rho^{n-4}U_{n-1}'(\rho m/2).
\]
Using \cref{lem:chebyshev} and $\rho\le1$ yields
\begin{equation}\label{eq:Bn}
 |\partial_sT_n|\le B_n:=
 \begin{cases}
 \displaystyle \frac{101}{200}(n-3)+\frac{11}{40}M_8 n,
   &n\text{ odd},\\[2mm]
 \displaystyle \frac{M_8n(101n-193)}{400},
   &n\text{ even}.
 \end{cases}
\end{equation}
The even bound retains the extra factor $c=\rho m/2$ in
\eqref{eq:U-parity}; losing that factor would make the estimate too weak.
By continuity the same bounds hold at $s=0$.
Now \cref{lem:coefficients} bounds the tail divided by $c_5$ by
\begin{equation}\label{eq:R8}
 R_8=\sum_{n\ge7}\frac{n-1}{4}\left(\frac5n\right)^8B_n.
\end{equation}
Exact summation through $n=100$ plus the integral bounds in
\cref{app:certificates} gives
\[
 R_8<0.388094585988858<\frac{39}{100}.
\]
The displayed decimal is only a readable outward bound; the replay
compares rational numbers exactly. The same majorant proves uniform
convergence of the differentiated tail.
\end{proof}

\begin{proof}[Proof of \cref{thm:main}]
At the zero, \cref{lem:low,lem:tail} give
\begin{equation}\label{eq:Ps-gap}
 P_s(m(s),s)>\left(\frac25-\frac{39}{100}\right)c_5
 =\frac{c_5}{100}>0.
\end{equation}
Implicit differentiation and \eqref{eq:Pm} imply
\[
 m'(s)=-\frac{P_s}{P_m}
 <-\frac{c_5}{120c_2}.
\]
Using \eqref{eq:normalization} proves \eqref{eq:main-bound}.
The equation at $s=0$ gives \eqref{eq:initial}; integration from zero
to $\rho$ proves \eqref{eq:gap}. Every estimate is uniform in $b>0$ and
in the full interval $0\le s\le1$.
\end{proof}

\begin{remark}[Why this is not a local Taylor argument]
Only finitely many modes are treated sharply, but the proof is not a
truncation without an error estimate. The inequality
$P_s>c_5/100$ incorporates \emph{every} mode $n\ge7$ and holds at
\emph{every} radius up to one. The Bernstein step retains the dependence
between $U,V_0,W$; replacing them by unrelated extreme values destroys
the useful positive margin.
\end{remark}

\section{Quantitative transfer from the infinite-inner-order family}
\label{sec:stability}
The limiting function is
$F_{a,\infty}(z)=\Li_a(z)-z$, whose strict global decrease is already
proved in the manuscript \cite{repo-cartesian}. Here we make that
limit stable with an explicit finite-$b$ criterion.

\subsection{A uniform perturbation lemma}
For $m,s,x,t\in[0,1]$ define
\begin{equation}\label{eq:f}
 f(m,s,x,t)=\frac{m-x-t}{D_x(m,s)D_t(m,s)}.
\end{equation}
The elementary bound
\begin{equation}\label{eq:D-bounds}
 \frac34\le D_u(m,s)\le2\qquad(0\le u,m,s\le1)
\end{equation}
follows by minimizing $u^2-mu$.

\begin{lemma}\label[lemma]{lem:perturbation}
On the unit four-dimensional cube,
\begin{equation}\label{eq:perturbation-constants}
 |f_t|\le\frac{304}{27}<12,
 \qquad |f_{tm}|,|f_{ts}|\le\frac{4224}{81}<64.
\end{equation}
Consequently, if $\epsilon=2^{-b}$ and
\begin{equation}\label{eq:Ginfty}
 G_\infty(m,s)=\E\frac{m-X}{D_X(m,s)},
\end{equation}
then, uniformly for $m,s\in[0,1]$,
\begin{equation}\label{eq:C1-stability}
 |G_b-G_\infty|\le12\epsilon,
 \quad |(G_b)_m-(G_\infty)_m|\le64\epsilon,
 \quad |(G_b)_s-(G_\infty)_s|\le64\epsilon.
\end{equation}
\end{lemma}
\begin{proof}
Put $Q=(D_xD_t)^{-1}$ and $N=m-x-t$. Then $|N|\le2$ and,
using \eqref{eq:D-bounds},
\[
 |Q|\le\frac{16}{9},\quad
 |Q_m|,|Q_s|,|Q_t|\le\frac{128}{27},\quad
 |Q_{tm}|\le\frac{1728}{81},\quad
 |Q_{ts}|\le\frac{1920}{81}.
\]
Here $|(D_t)_t|\le2$, $|(D_u)_m|,|(D_u)_s|\le1$,
$|(D_t)_{tm}|\le1$, and $|(D_t)_{ts}|\le2$ give the bounds directly
by the product rule. Thus
\[
 |f_t|\le |Q|+2|Q_t|\le\frac{304}{27},
\]
\[
 |f_{tm}|\le |Q_m|+|Q_t|+2|Q_{tm}|
 \le\frac{4224}{81},\qquad
 |f_{ts}|\le |Q_s|+2|Q_{ts}|\le\frac{4224}{81}.
\]
The fundamental theorem of calculus in $t$ compares $t=XV$ with
$t=0$. Since $0\le XV\le V$ and $\E V=2^{-b}$, expectation yields
\eqref{eq:C1-stability}. The bounded derivatives also justify passing
$m,s$ derivatives through the expectation.
\end{proof}

\subsection{A variance gap that is uniform in radius}
\begin{lemma}\label[lemma]{lem:variance}
For $0\le m,s\le1$,
\begin{equation}\label{eq:limiting-derivatives}
 (G_\infty)_m=\E D_X^{-2}\ge\frac14,
 \qquad
 (G_\infty)_s=\E\frac{X(X-m)^2}{D_X^2}
 \ge\frac{\Delta(a)}4.
\end{equation}
Also
\begin{equation}\label{eq:endpoint-infty}
 G_\infty(1,s)\ge\frac{1-(2/3)^a}{2},
 \qquad 0<\Delta(a)<1-(2/3)^a<1.
\end{equation}
\end{lemma}
\begin{proof}
Differentiation of \eqref{eq:Ginfty} gives the two identities in
\eqref{eq:limiting-derivatives}. By $D_X\le2$ and completing the
square in $m$,
\begin{align*}
 (G_\infty)_s
 &\ge\frac14\E[X(X-m)^2]\\
 &\ge\frac14\left(\E X^3-\frac{(\E X^2)^2}{\E X}\right)
 =\frac14\left((2/5)^a-(3/8)^a\right).
\end{align*}
The difference is strictly positive because $2/5>3/8$.
At $m=1$, $1-X\ge0$ and $D_X\le2$, proving the endpoint bound.
Finally
\[
 \Delta(a)=(2/5)^a[1-(15/16)^a]<1-(2/3)^a.
\]
\end{proof}

\begin{proof}[Proof of \cref{thm:large-b}]
Assume \eqref{eq:large-b-condition}. Since $\Delta(a)<1$, this implies
$b>8$ and hence $a+b>1$, so the inherited uniqueness theorem
\eqref{eq:inherited-angle} applies. Equations
\eqref{eq:C1-stability} and \eqref{eq:limiting-derivatives} give
\[
 (G_b)_m>\frac14-\frac{\Delta(a)}4>0,
 \qquad
 (G_b)_s>\frac{\Delta(a)}4-\frac{\Delta(a)}4=0.
\]
Moreover $G_b(0,s)<0$ directly from \eqref{eq:G-b}, whereas
\[
 G_b(1,s)>\frac{1-(2/3)^a}{2}-\frac{3\Delta(a)}{64}>0.
\]
There is therefore a simple root in $0<m<1$, and it is the unique
upper-semicircle root. Its derivative is
$m'(s)=-(G_b)_s/(G_b)_m<0$, which proves the asserted direction.

Under $\epsilon\le\Delta(a)/512$, the same estimates give
$(G_b)_s\ge\Delta(a)/8$ and
\[
 (G_b)_m\le\frac{16}{9}+\frac{\Delta(a)}8<2.
\]
Thus $m'(s)\le-\Delta(a)/16$. Equation
\eqref{eq:normalization} proves \eqref{eq:large-b-gap}.
\end{proof}

\begin{corollary}[A compact outer-order interval]\label[corollary]{cor:compact}
Fix $0<\alpha\le A<\infty$. If
\[
 2^{-b}<\frac1{256}\left(\frac25\right)^A
                  \left[1-\left(\frac{15}{16}\right)^\alpha\right],
\]
then strict global decrease holds simultaneously for every
$a\in[\alpha,A]$.
\end{corollary}
\begin{proof}
The displayed numerator is a lower bound for
$\Delta(a)=(2/5)^a[1-(15/16)^a]$ throughout that interval.
\end{proof}

The perturbation theorem supplies more than pointwise convergence as
$b\to\infty$. It controls the zero equation and both of its relevant
derivatives uniformly in radius, and gives a positive lower gap before
the perturbation is made. Without these two properties, a strict sign
in the limit need not persist on the entire interval.

\section{An exact cubic-moment criterion for radial motion}
\label{sec:cubic}
The previous arguments prove two sign regimes. The following identity
isolates the obstruction in the remaining cases. It is stated on
$a,b>0$, $a+b>1$, where \eqref{eq:inherited-angle} supplies the required
simple root and sign. The algebra also holds at any other simple root
satisfying the same geometric hypotheses.

Write $\widehat\eta(s)=\eta_{a,b}(\sqrt s)=m(s)/2$, and, at its zero,
put
\begin{equation}\label{eq:cubic-defs}
 A=1-s\widehat\eta^2>0,\qquad
 D=1+s(u^2-2\widehat\eta u),\qquad
 M_j=\int (u-\widehat\eta)^jD^{-3}\dd\mu(u).
\end{equation}
The inequality $A>0$ follows from $0<\widehat\eta<1$ and $s\le1$.

\begin{theorem}[Weighted cubic-moment velocity]\label{thm:cubic}
At the angular zero,
\begin{equation}\label{eq:root-moment}
 A M_1+sM_3=0
\end{equation}
and
\begin{equation}\label{eq:velocity}
 \widehat\eta'(s)=
 -\frac{2M_3}{A^2M_0-3sAM_2+4s^2\widehat\eta M_3}.
\end{equation}
The denominator is strictly positive. Consequently the sign of
$\eta_{a,b}'(\rho)$ is the opposite of the sign of $M_3$.
\end{theorem}
\begin{proof}
Set
\[
 \mathcal G(\eta,s)=\int(\eta-u)
           [1+s(u^2-2\eta u)]^{-2}\dd\mu(u).
\]
Equation \eqref{eq:P-integral} gives $P(2\eta,s)=2\mathcal G(\eta,s)$.
Since $D=A+s(u-\eta)^2$, the equation $\mathcal G=0$ is exactly
\eqref{eq:root-moment}. Holding $\eta$ fixed when differentiating in
$s$,
\[
 \mathcal G_s=2(M_3-\eta^2M_1)=\frac{2M_3}{A}.
\]
Similarly,
\[
 \mathcal G_\eta=A M_0-3sM_2-4s\eta M_1
 =A M_0-3sM_2+\frac{4s^2\eta M_3}{A}.
\]
Implicit differentiation gives \eqref{eq:velocity}.
At the inherited angular zero, the angular derivative of the imaginary
part is negative, while $\partial_\theta m=-2\sin\theta/\rho<0$.
Equation \eqref{eq:P-integral} therefore gives $P_m>0$.
As $\mathcal G_\eta=P_m(2\eta,s)>0$, the denominator in
\eqref{eq:velocity} is $A\mathcal G_\eta>0$.
Finally $\eta'(\rho)=2\rho\widehat\eta'(s)$, which does not change sign.
\end{proof}

\begin{corollary}[A precise reformulation of the remaining conjecture]
For integer $a\ge2$, global normalized-radius decrease is equivalent
to positivity of
\[
 \int_0^1
 \frac{(u-\eta_{a,b}(\rho))^3}
 {[1+\rho^2(u^2-2\eta_{a,b}(\rho)u)]^3}\dd\mu_{a,b}(u)
\]
at every $0<\rho\le1$. The integral is strictly positive in both
regimes of \cref{thm:main,thm:large-b}.
\end{corollary}
This moment is centered at the root parameter, not necessarily at the
mean of the reweighted measure. Calling it a skewness criterion is
useful shorthand, but it is not the usual normalized statistical
skewness. Positivity of $\mu$ alone does not determine its sign.

\section{Positive binomial identities and exact remainders}
\label{sec:identities}
Let $g_{a,b}=\Imn F_{a,b}(\ii)$. The squared-resolvent representation
makes a particularly simple family of alternating finite sums positive.

\begin{theorem}[Positive Gaussian binomial transform]\label{thm:binomial}
For $a,b>0$ and every integer $k\ge0$, define
\begin{equation}\label{eq:Bk}
 \mathcal B_k(a,b)=\sum_{j=0}^k(-1)^j\binom kj
   \frac{H_{2j+2}^{(b)}}{(2j+2)(2j+3)^a}.
\end{equation}
Then $\mathcal B_k(a,b)>0$, and
\begin{equation}\label{eq:binomial}
 -g_{a,b}=\sum_{k=0}^\infty\frac{k+1}{2^{k+1}}\mathcal B_k(a,b).
\end{equation}
For the partial sum $S_N$ over $0\le k<N$, the remainder is exactly
\begin{equation}\label{eq:binomial-remainder}
 -g_{a,b}-S_N
 =2^{-N}\int_0^1
 \frac{u(1-u^2)^N[(N+2)+Nu^2]}{(1+u^2)^2}\dd\mu(u),
\end{equation}
and hence, for every $N\ge1$,
\begin{equation}\label{eq:binomial-error}
 0<-g_{a,b}-S_N<\frac{N+2}{2^{N+1}}\,
             \frac{H_2^{(b)}}{3^a}.
\end{equation}
\end{theorem}
\begin{proof}
By \eqref{eq:squared},
\[
 \mathcal B_k(a,b)=\int u(1-u^2)^k\dd\mu(u)>0,
 \qquad -g_{a,b}=2\int\frac{u}{(1+u^2)^2}\dd\mu(u).
\]
Set $q=(1-u^2)/2\in[0,1/2]$. The identity
$(1-q)^{-2}=\sum_{k\ge0}(k+1)q^k$ gives \eqref{eq:binomial}; all
terms are nonnegative, so monotone convergence justifies integration.
The finite algebraic remainder
\[
 \sum_{k=N}^\infty(k+1)q^k
 =q^N\frac{(N+1)-Nq}{(1-q)^2}
\]
gives \eqref{eq:binomial-remainder} after substitution.
For $v\in[0,1]$,
$[(N+2)+Nv]/(1+v)^2\le N+2$. Since $\int u\dd\mu=c_3/2$,
\eqref{eq:binomial-error} follows, with strictness because $\mu$
has positive mass in the open interval.
\end{proof}

This theorem concerns a different transformation from the manuscript's
Euler transform of odd harmonic coefficients. The one-sided remainder
and the positivity of each $\mathcal B_k$ are useful structural facts;
no claim is made that the bound $O(N2^{-N})$ supersedes the sharper
existing $O(2^{-N})$ evaluator \cite{incoming-euler}.

\subsection{Three explicit specializations}
\begin{corollary}[An identity for $\pi\log2$]\label[corollary]{cor:pi-log}
\begin{equation}\label{eq:pi-log}
 \pi\log2=4\sum_{k=0}^\infty\frac{k+1}{2^k}
 \sum_{j=0}^k(-1)^j\binom kj
      \frac{H_{2j+2}}{(2j+2)(2j+3)}.
\end{equation}
Every inner finite sum is positive, and the error after $N$ outer terms
is strictly between zero and $(N+2)2^{1-N}$.
\end{corollary}
\begin{proof}
The classical identity $F_{1,1}(z)=\tfrac12\log^2(1-z)$ follows directly
by differentiating its series and matching the value at zero.
At $z=\ii$, $\log(1-\ii)=\tfrac12\log2-\ii\pi/4$, whence
$-g_{1,1}=\pi\log2/8$. Apply \cref{thm:binomial} with
$c_3=H_2/3=1/2$ and multiply the value and error bound by eight.
\end{proof}

Define $\beta(t)=\Imn\Li_t(\ii)$, the Dirichlet beta function with
analytic continuation when necessary. Taking the endpoint limits in
\cref{thm:binomial} gives two further identities.
\begin{corollary}[Beta differences]\label[corollary]{cor:beta}
For $a>0$,
\begin{align}\label{eq:beta-zero}
 \beta(a)-\beta(a-1)
 &=\frac12\sum_{k=0}^\infty\frac{k+1}{2^k}
       \sum_{j=0}^k(-1)^j\binom kj(2j+3)^{-a},\\
 1-\beta(a)
 &=\frac12\sum_{k=0}^\infty\frac{k+1}{2^k}
       \sum_{j=0}^k(-1)^j\binom kj
                   \frac1{(2j+2)(2j+3)^a}.
 \label{eq:beta-infty}
\end{align}
Every inner sum in these formulas is strictly positive.
In particular $a=2$ in \eqref{eq:beta-zero} gives an explicit positive
series for Catalan's constant minus $\pi/4$.
\end{corollary}
\begin{proof}
As $b\downarrow0$, $V\to1$ in probability in
\eqref{eq:probability}; as $b\to\infty$, $V\to0$ in probability.
All relevant kernels at $z=\ii$ are bounded and continuous. The limits
are respectively
$F_{a,0}(z)=\Li_{a-1}(z)-\Li_a(z)$ and
$F_{a,\infty}(z)=\Li_a(z)-z$.
Their positive squared-resolvent measures are, respectively, $c_2$
times the law of $X$ and $c_2$ times the mixture of the uniform laws on
$[0,X]$. Both have full support. Their moments are the limits of
\eqref{eq:squared}, which give the two displayed inner sums.
Apply the same proof as \cref{thm:binomial} to these limiting measures.
This avoids relying on an unjustified interchange of two infinite sums.
\end{proof}

\subsection{The generating function of the positive moments}
For $|t|<1$, another useful exact identity is
\begin{equation}\label{eq:B-generating}
 \sum_{k=0}^\infty\mathcal B_k(a,b)t^k
 =\int_0^1\frac{u}{1-t+tu^2}\dd\mu(u).
\end{equation}
The geometric expansion is uniformly convergent on compact subdisks.
For $0<t<1$ the right-hand side is positive. More precisely,
$\left.\tfrac12\frac{d}{dt}(t\sum_k\mathcal B_kt^k)\right|_{t=1/2}
=-g_{a,b}$, which recovers \eqref{eq:binomial}. More generally, all derivatives in $t\in[0,1)$
are positive, providing a generating-function form of the finite-sum
positivity.

\section{Reproducible exact computation}
\label{sec:computation}
All acceptance decisions in the supplied verifiers use Python's
standard-library \path{fractions.Fraction}. The analytic arguments,
not finite sampling, transfer those decisions to statements about the
infinite functions and continuous parameter ranges.

\subsection{Global certificate}
The verifier \path{code/verify_global.py} checks the ten coefficients
in \cref{tab:bernstein}, the root-strip and radial-tail inequalities,
the rational comparisons extending the core to real $a\ge8$, and the
eight integer thresholds for \cref{thm:large-b}.
The finite sums use cutoff $100$. Representative exact outward bounds
printed by the verifier are
\begin{center}
\begin{tabular}{@{}ll@{}}
\toprule
Certified quantity & Outward decimal information\\
\midrule
Upper bound for $B_8$ & $<0.302408858414858$\\
Localization margin & $>0.022510694406543$\\
Upper bound for $R_8$ & $<0.388094585988858$\\
Low-mode contribution divided by $c_5$ & $>2/5$\\
Full $P_s/c_5$ at the root & $>1/100$\\
\bottomrule
\end{tabular}
\end{center}
These decimals are not asserted to be optimal constants.

\subsection{Independent angular-root certificates}
There is an exponentially convergent rational evaluator for
\eqref{eq:P-integral}. On $0\le m\le1/3$, $0\le s\le1$, define
\begin{equation}\label{eq:E-eval}
 E(u)=1-\frac23D_u=\frac13+\frac{2sm}{3}u-\frac{2s}{3}u^2.
\end{equation}
Since $35/36\le D_u\le2$,
$|E(u)|\le19/54<3/8$. Thus
\begin{equation}\label{eq:P-eval}
 P_N(m,s)=\frac49\sum_{k=0}^{N-1}(k+1)
             \int(m-2u)E(u)^k\dd\mu(u)
\end{equation}
is computable from finitely many rational moments when $a,b,m,s$ are
rational with integer $a,b$. Its error is bounded by
\begin{equation}\label{eq:P-error}
 |P-P_N|\le\frac89 c_2
       \left(\frac38\right)^N
       \frac{(N+1)-3N/8}{(1-3/8)^2}.
\end{equation}
Indeed $|m-2u|\le2$, and the differentiated geometric tail used in
\cref{thm:binomial} bounds the absolute series. This estimate is valid
even though $E$ can be negative.

The file \path{data/root_brackets.json} records candidate rational
endpoints. Their method of discovery has no evidentiary role: the
verifier accepts a bracket only when the upper error bound for $P$ at
the lower endpoint is negative and the lower error bound at the upper
endpoint is positive. Monotonicity in $m$ then certifies the unique
root. The replay uses $N=36$ and certifies the following enclosures for
$\eta=m/2$; the displayed endpoints are rounded outward.
\begin{center}
\begin{tabular}{@{}rrrrr@{}}
\toprule
$a$&$b$&$\rho$&Lower bound for $\eta$&Upper bound for $\eta$\\
\midrule
8&1&$1/2$&0.029174544600&0.029174544751\\
8&1&$1$  &0.028924151300&0.028924151451\\
8&2&$1/2$&0.024321334850&0.024321335001\\
8&2&$1$  &0.024137504900&0.024137505051\\
9&1&$1/2$&0.019469017900&0.019469018051\\
9&1&$1$  &0.019354467350&0.019354467501\\
12&1&$1/2$&0.005777259000&0.005777259151\\
12&1&$1$  &0.005767707400&0.005767707551\\
\bottomrule
\end{tabular}
\end{center}
These checks illustrate and independently test the normalization and
root evaluator. They are not used as a substitute for the uniform proof.

\subsection{Gaussian enclosures and an elementary control}
The script \path{code/verify_identities.py} computes 80 terms of
\eqref{eq:binomial} for nine integer parameter pairs. It verifies 720
strictly positive rational values of $\mathcal B_k$, and separately
checks 100 polynomial geometric-remainder identities. The following
are exact outward enclosures for $-g_{a,b}$:
\begin{center}
\small
\begin{tabular}{@{}rrll@{}}
\toprule
$a$&$b$&Lower bound&Upper bound\\
\midrule
1&1&0.272198261287950266&0.272198261287950267\\
1&2&0.251642673723681485&0.251642673723681486\\
2&1&0.113648917599903643&0.113648917599903644\\
2&2&0.101022920883400012&0.101022920883400013\\
3&4&0.032901765615888080&0.032901765615888081\\
8&1&0.000223662993166255&0.000223662993166256\\
8&2&0.000187104280701129&0.000187104280701130\\
9&4&0.000053452648646007&0.000053452648646008\\
12&1&0.000002814149505465&0.000002814149505466\\
\bottomrule
\end{tabular}
\end{center}
For $(a,b)=(1,1)$ an independent interval for $\pi\log2/8$, obtained
from rational arctangent and logarithm series, lies strictly inside
the unrounded binomial enclosure. This checks the scale and sign
against an elementary special value without relying on a floating-point
polylogarithm implementation. Details are in \cref{app:control}.

\section{Source audit, corrections, and integration boundaries}
\label{sec:audit}
\subsection{What should change in the manuscript}
The final remark of the local-radius section should no longer leave
\emph{all} integer outer orders in the global conjectural column.
It can say that full-radius decrease is proved for every real $a\ge8$,
and for every $a>0$ under \eqref{eq:large-b-condition}, while preserving
the remaining cases as open. The left-endpoint limit theorem in the
Cartesian-motion section should be followed by the quantitative
finite-$b$ transfer rather than being replaced or presented as new.

The new binomial identities can be added alongside positive-kernel
identities. Their computational status should remain distinct from
the sharp Euler theorem. The supplied \path{integration/} files use
namespaced labels and provide an additive insertion, not an automatic
rewrite of the canonical manuscript.

\subsection{A small existing notation correction}
The proof in \path{05-real-positive-kernel.tex} of the
Hurwitz--polylogarithm transport theorem contains the phrases ``with $q=n$'' and
``boundedness of $h$''. These are inconsistent with the remainder kernel
notation $q(t)$ used in the same formulas. The intended substitution
is the positive scalar argument of the Hurwitz formula equal to $n$;
the bounded kernel in the displayed integral is $q(t)$.
A safe wording is: ``Apply the Hurwitz difference formula at the
positive integer argument $n$; use boundedness of $q(t)$ near zero and
its bounded large-$t$ behavior.'' This is a notation repair, not a
refutation of the transport identity. It should be checked against the
surrounding notation at integration time. These editorial repairs should be deduplicated against incoming audit
notes during integration; no claim of first discovery is made.

\subsection{Non-errors that must not become errors during integration}
The inspected local theorem is correctly restricted to a neighborhood
of the origin. We found no new fatal flaw in that proof. A local
coefficient sign must not be cited as a global radial sign. Likewise,
the Cartesian quantity $\rho\cos\theta$ and the normalized quantity
$\cos\theta/\rho$ have different derivatives: the manuscript's
outward-abscissa theorem does not settle the normalized conjecture.

The incoming fractional bifurcation article has a global statement in
the inner-order parameter \emph{along a quadratic threshold}, but its
turning-point geometry is local in radius. Those quantifiers should
not be combined with the all-radius quantifiers of the present paper.
The new large-$a$ theorem does not erase the fractional counterexamples.

Finally, the positive squared-resolvent measure is not the manuscript's
signed first-resolvent measure, and analytic positivity is not an
arithmetic depth-reduction certificate. The $S_6$ and surviving $S_8$
candidates keep their existing status. We do not present earlier
all-depth oscillation, the sharp Euler constant, or the local quartic
transition as newly solved here.

\section{Further research questions and proposed directions}
\label{sec:future}
\subsection{The six remaining integer strips}
The most concrete next target is
\begin{conjecture}[Remaining integer normalized-radius problem]
For $a\in\{2,3,4,5,6,7\}$ and every $b>0$,
$\eta_{a,b}'(\rho)<0$ on $0<\rho\le1$.
\end{conjecture}
By \cref{thm:large-b}, only $0<b<B_a$ from the threshold table remains.
The squared-resolvent evaluator is uniform in radius and avoids slowly
converging unit-circle derivatives. This suggests a genuine
computer-assisted proof strategy: cover each bounded $(b,s)$ region,
isolate the root in $m$, and certify positivity of $P_s$ there using
interval moments and analytic tails. Such a proof would require
outward bounds for the noninteger powers in $H_n^{(b)}$, and a treatment
of the edges $b\downarrow0$ and $s\downarrow0$. No such complete
parameter-covering certificate is included here.

\subsection{The first outer order}
The manuscript predicts decrease for $a=1,b>1$, increase for
$a=1,0<b<1$, and the constant $\eta=1/2$ at $(1,1)$.
Our transfer theorem proves the decreasing claim only for $b\ge14$.
The range $1<b<14$ remains a focused target. An integral comparison
against the exact logarithmic solution may be more efficient there
than a mode-by-mode estimate.

\subsection{How small can the uniform outer threshold be?}
The number eight is an explicit sufficient threshold of this proof,
not a demonstrated phase boundary. Retaining favorable signs from
more even and odd modes, or certifying the core in a smaller actual
parameter region $y=x^{\log_2 3}$ rather than $x^2\le y\le x$, may lower
it. A substantially different question asks for the least real $A$
such that global decrease holds for all $a\ge A,b>0$. Fractional
turning-point phenomena make it inappropriate to infer the answer
from the local $a\ge2$ result alone.

\subsection{Sharp stability in the inner order}
The constants 12, 64, and 256 arise from deliberately uniform kernel
bounds. Determine the least sufficient $B(a)$, or an asymptotically
sharp $B(a)$ as $a\downarrow0$ and $a\to\infty$. The proof already
isolates the useful variance gap
$\Delta(a)=\E X^3-(\E X^2)^2/\E X$; retaining $\E(XV)$ rather than
bounding it by $\E V$, and exploiting the root equation, should improve
the perturbation threshold. These observations are proposals, not
claims of a sharp asymptotic formula.

\subsection{A structural cubic-moment theorem}
The criterion in \cref{thm:cubic} reduces radial motion to a sign of a
reweighted third moment. Characterize positive measures $\mu$ for
which that sign is positive at every root. A useful theorem should
state additional shape or transform hypotheses, because bare
positivity cannot suffice. The harmonic moment family offers an
explicit test class with known fractional changes of sign.

\subsection{Further identities and arithmetic certificates}
Equation \eqref{eq:B-generating} suggests specializations at rational
$t$, derivatives of arbitrary order, and cyclotomic analogues of the
positive Gaussian transform. Their analytic convergence and remainder
bounds are accessible from the same measure. The separate question
is whether any such identity yields a finite, meaningful reduction
in the manuscript's arithmetic period basis. The $S_6,S_8$ targets
require that second step; an accelerated positive infinite series is
not itself the missing finite arithmetic certificate.

\subsection{Deeper strict polylogarithms}
The incoming all-depth report controls several angular roots for
integer leading indices. It is natural to ask which normalized
Cartesian displacements of those roots have uniform radial signs,
and whether positive higher-resolvent measures lead to analogous
central-moment identities. Multiple-star sums and strict sums must
remain separate, as must the number of angular roots and the
monotonicity of an individual normalized branch.

\appendix
\section{Finite certificates and analytic tail bookkeeping}
\label{app:certificates}
Let $N=100$ and $M=M_8$. The following are rational upper bounds for
the sums in \eqref{eq:certificate-three}:
\begin{align*}
 \overline B&=\sum_{n=3}^{N}n(n-1)(2/n)^8
                 +\frac{2^8}{5N^5},\\
 \overline O&=\sum_{\substack{3\le n\le N\\n\text{ odd}}}
                    \frac{n-1}{2}(3/n)^8
                 +\frac{3^8}{12N^6},\\
 \overline E&=\sum_{\substack{4\le n\le N\\n\text{ even}}}
                    n(n-1)(2/n)^8
                 +\frac{2^8}{5N^5}.
\end{align*}
For example, $n(n-1)n^{-8}\le n^{-6}$ and
$\sum_{n>N}n^{-6}\le\int_N^\infty t^{-6}\dd t=1/(5N^5)$.
For the odd sum, use $(n-1)n^{-8}/2\le n^{-7}/2$ and bound the
remaining odd indices by \emph{all} indices; this is deliberately
conservative. The same method bounds the radial sum by
\begin{align}\label{eq:exact-tail-upper}
 \overline R={}&\sum_{n=7}^{N}\frac{n-1}{4}(5/n)^8 B_n\notag\\
 &+\frac{5^8}{4}\left(\frac{101}{200}+\frac{11M}{40}\right)
                  \frac1{5N^5}
 +\frac{5^8\,101M}{1600}\frac1{4N^4}.
\end{align}
Indeed the odd summand is bounded by
$5^8(101/200+11M/40)n^{-6}/4$, and the even summand by
$5^8\,101M n^{-5}/1600$.

The verifier performs the exact rational comparisons
\[
 \overline B<\frac13,\quad
 C\left(1-\frac{101}{200}\overline E\right)
       -\frac{101}{100}\overline O>\frac1{50},\quad
 \overline R<\frac{39}{100}.
\]
For a polynomial $p(x)=\sum_{k=0}^d p_kx^k$, the degree-$d$ Bernstein
coefficients used in \cref{tab:bernstein} are
\[
 b_j=\sum_{k=0}^j p_k\frac{\binom jk}{\binom dk}.
\]
The code computes these from \eqref{eq:core-polynomial} over
$\mathbb Q$, compares them with the displayed integers divided by
the common denominator, and checks that every coefficient is positive.
The Bernstein basis functions are nonnegative and sum to one on
$[0,1]$, so these finite signs imply positivity on the full interval.

There is no sampling in $a$ or $b$ in this proof: the exponential
comparison proves the extension from $a=8$ to all real $a\ge8$, and
the affine/Bernstein argument covers every $b>0$ through
$x^2\le y\le x$. Only the elementary infinite sums are truncated,
with the rigorous integral tails displayed above.

\section{Independent control of the elementary special value}
\label{app:control}
For rational $0<x<1$, consecutive partial sums of
\[
 \arctan x=\sum_{k=0}^\infty\frac{(-1)^kx^{2k+1}}{2k+1}
\]
bracket the exact value by the alternating-series theorem.
The verifier uses 100 terms at $x=1/5$ and 35 at $x=1/239$, together
with Machin's identity
\[
 \pi=16\arctan(1/5)-4\arctan(1/239).
\]
For a direct check of the identity, the tangent addition formula gives
$\tan(4\arctan(1/5))=120/119$ and
$\tan(4\arctan(1/5)-\arctan(1/239))=1$; the angle is in $(0,\pi/2)$.

For $\log2$, use
\[
 \log2=2\sum_{k=0}^\infty\frac{(1/3)^{2k+1}}{2k+1}.
\]
After $L$ terms the positive remainder is at most
\[
 \frac{2(1/3)^{2L+1}}{(2L+1)(1-1/9)}.
\]
The verifier takes $L=110$, multiplies the positive intervals for
$\pi$ and $\log2$, divides by eight, and verifies strict containment
inside the rational $(1,1)$ Gaussian enclosure. Every endpoint in
this independent calculation is rational.

\section{Reproduction, dependency map, and integration}
\label{app:reproduce}
From the package root, with Python 3.10 or newer and without optimized
mode, run
\begin{verbatim}
python code/run_all.py
\end{verbatim}
The individual scripts are
\begin{verbatim}
python code/verify_global.py
python code/verify_roots.py
python code/verify_identities.py
\end{verbatim}
They regenerate their receipts under \path{data/}. Run them on a copy
when preservation of delivered evidence bytes matters. No third-party
Python dependency is needed. To rebuild the PDF with a standard TeX
installation, run \path{build.sh}, or run \path{pdflatex article.tex}
twice. The bibliography is included in the source, so no bibliography
processor or network access is required.

The proof dependencies are deliberately limited. The $a\ge8$ theorem
uses the inherited positive kernel but proves its own root localization
and uniqueness. The finite-$b$ transfer uses the inherited uniqueness
result only to identify the constructed $0<m<1$ root as the sole upper
angular zero. The cubic criterion uses the inherited angular-derivative
sign. The binomial identities depend only on the positive squared
resolvent, not on the radial theorem. The exact root and Gaussian
brackets test concrete values and normalizations independently of
the global parameter argument.

The package is additive and makes no remote repository changes.
Suggested placement is a new report directory under
\path{Analysis/Polylogarithms/docs/reports/}, with the namespaced excerpt
from \path{integration/} inserted in the real-order chapter after the
local-radius material. The status ledger should record the two proved
global regimes and preserve the unresolved strips, the first-order
cases not covered here, and the arithmetic conjectures separately.

\begin{thebibliography}{20}
\small\raggedright
\bibitem{repo-local}
ProveIt Contributors,
\emph{A proved local part of the normalized-radius conjecture},
\path{Analysis/Polylogarithms/docs/manuscript/chapters/05-real-local-radius.tex},
Git blob \texttt{30e3e31944be5bc0d03a8485386dd1608c075507}.
Repository \url{https://github.com/VladimirReshetnikov/ProveIt},
accessed 10 October 2026.

\bibitem{repo-cartesian}
ProveIt Contributors,
\emph{Cartesian motion of the complete real-order zero arc},
\path{Analysis/Polylogarithms/docs/manuscript/chapters/05-real-cartesian-motion.tex},
Git blob \texttt{0deb9bb54adaef8783d4e0e7ce43d426643450a0}.
Same repository and access date.

\bibitem{repo-kernel}
ProveIt Contributors,
\emph{A positive two-variable representation at all positive orders},
\path{Analysis/Polylogarithms/docs/manuscript/chapters/05-real-positive-kernel.tex},
Git blob \texttt{044d3825b90dce437ac9007733e6447cf9563098}.
Same repository and access date.

\bibitem{repo-orders}
ProveIt Contributors,
\emph{Real orders: signed-measure domain and consequences for angular zeros},
\path{Analysis/Polylogarithms/docs/manuscript/chapters/05-real-orders.tex},
Git blob \texttt{61db4b7721bfdc78434c1efc9e932e15e1c79f33}.
Read at revision
\texttt{3338e535ef54ad8f6be56bdce8b3d084a2924465}.

\bibitem{incoming-bifurcation}
ProveIt research continuation,
\emph{A Unique Quartic Transition, M\"obius Rigidity, and Two Radial
Turning Points for Real-Order Double Polylogarithms},
10 October 2026. Accessible delivered source
\path{polylogarithms_radial_bifurcation.tex}; corresponding incoming
archive \path{ProveIt_Polylogarithms_Radial_Bifurcation_2026-10-10.zip}.
Scope and contribution statements consulted; not a full independent
replay of that report's certificates.

\bibitem{incoming-oscillation}
ProveIt research continuation,
\emph{Signed-Kernel Oscillation at Every Depth: Minimal Pick Compensation,
Zero Geometry, and Certified Polylogarithmic Identities},
10 October 2026. Corresponding incoming archive
\path{ProveIt_AllDepth_Oscillation.zip}. Abstract and scope/audit excerpts
consulted; not a full proof audit.

\bibitem{incoming-euler}
ProveIt research continuation,
\emph{The Sharp Universal Euler Constant for Harmonic Polylogarithms},
10 October 2026. Corresponding incoming archive
\path{ProveIt_Sharp_Universal_Euler_2026-10-10.zip}.
Abstract and contribution statements consulted.

\bibitem{dlmf-polylog}
NIST Digital Library of Mathematical Functions,
\S25.12, \emph{Polylogarithms}, in particular the defining series and
Mellin integral, equations 25.12.10--11.
\url{https://dlmf.nist.gov/25.12}, accessed 10 October 2026.

\bibitem{dlmf-chebyshev}
NIST Digital Library of Mathematical Functions,
\S18.5, \emph{Explicit Representations}, trigonometric representations
of the Chebyshev polynomials.
\url{https://dlmf.nist.gov/18.5}, accessed 10 October 2026.
\end{thebibliography}
\end{document}
